% Context from chapters-tex/approximation.tex; for mathematical reference, not compilation.
\section{Fourier Linear Approximation of Smooth Functions}

For integer $\alpha$, the smooth model~\eqref{eq-smooth-model} bounds all continuous derivatives through that order. Greater regularity of $f$, measured by a larger $\al$, permits faster approximation. Figure~\ref{fig-signal-model-sobolev} illustrates the related Sobolev measure of smoothness.

                                        
\subsection{1-D Fourier Approximation}

A 1-D signal $f \in \Ldeux([0,1])$ is associated with a 1-periodic function $f(t+1)=f(t)$ defined for $t \in \RR$.

  
\paragraph{Low pass approximation.}

For even $M$, consider the linear Fourier approximation that retains frequencies up to $M/2$:
\eq{
	f_M^{\text{lin}} = \sum_{m=-M/2}^{M/2} \dotp{f}{e_m}e_m
}
where we use the 1-D Fourier atoms
\eq{
	\foralls m\in\ZZ,\quad e_m(t) \eqdef e^{2 i \pi m t}.
}
This symmetric cutoff retains $M+1$ Fourier atoms; the extra coefficient does not affect the asymptotic rates below.

Figure \ref{fig-ourier-approx-1d} shows these approximations for increasing $M$. The jumps in $f$ produce substantial error and ringing near the singularities.

                                                                                                                                                               

\myfigure{
\tabdeux{
\image{approximation}{.48}{fourier-approx-1d-0}&
\image{approximation}{.48}{fourier-approx-1d-1}\\
\image{approximation}{.48}{fourier-approx-1d-2}&
\image{approximation}{.48}{fourier-approx-1d-3}
}
}{ 
	Fourier approximation of a signal.  
}{fig-ourier-approx-1d}

                                       
                                                                                                                                                                                                      

                                          
                                                                                                                                                                   


This low-pass approximation is a convolution, since
\eq{
	f_M = \sum_{m=-M/2}^{M/2} \dotp{f}{e_m} e_m = f \star h_M
	\qwhereq
	\hat h_{M} \eqdef 1_{[-M/2,M/2]}.
}
The kernel $h_M$ is the Dirichlet kernel introduced in the Fourier chapter.

The following elementary bound shows that this approximation error decays for $\Cc^\al$ signals.

\begin{prop}
	For integer $\alpha\geq1$ and $f\in C^\alpha(\RR/\ZZ)$
	\eq{
		\norm{ f-f_M^{\text{lin}} }^2 = O(M^{-2\al+1}).
	}
\end{prop}

\begin{proof}
Using integration by parts, one shows that for a $\Cal$ function $f$ in the smooth signal model~\eqref{eq-smooth-model},
$\Ff(f^{(\ell)})_m = (2\imath m \pi)^\ell \hat f_m$, so that 
Cauchy--Schwarz gives
\eq{
	|\dotp{f}{e_m}| \leq |2\pi m|^{-\al} \norm{f^{(\al)}}_2.
}
This implies
\eq{
	\norm{ f-f_M^{\text{lin}} }^2 = \sum_{|m|>M/2} |\dotp{f}{e_m}|^2
	\leq \norm{f^{(\al)}}^2_2 \sum_{|m|>M/2} |2\pi m|^{-2\al} = O(M^{-2\al+1}).
}
\end{proof}

Using the full energy of $f^{(\al)}$ gives a sharper bound under the same assumptions. The argument also applies to the larger Sobolev class.

\begin{prop}\label{prop-sobol-fourier-lin}
	For $\alpha>0$ and a periodic signal in the Sobolev model~\eqref{eq-sobolev-model}, one has
	\eql{\label{eq-sovolev-signal-decay}
		\norm{f-f^{\text{lin}}_M}^2 \leq C \norm{f}_{\mathrm{Sob}(\alpha)}^2  M^{-2\al}.
	}
\end{prop}
\begin{proof}
	One has
	\eqml{
		\norm{f}_{\mathrm{Sob}(\alpha)}^2 &= \sum_m |2 \pi m|^{2\al} |\dotp{f}{e_m}|^2
		\geq \sum_{|m|>M/2} |2 \pi m|^{2\al} |\dotp{f}{e_m}|^2 \\ 
		& \geq (\pi M)^{2\al} 